% Additions for the simulated paper. main.tex reads this file only when \simulatedpaper is set.
% \simnum{...} marks a PROJECTED number (not measured). It prints the number unchanged, or in red with \showsimulated.
% List every projected number:  grep -n -o 'simnum{[^}]*}' sim/*.tex
%
% How each projection was made (2026-11-13). Replace them with measurements as the runs finish.
%  Table 1, InternVL3.5-4B rows Top-K (w/ pruning), Q-Frame, TALON: additive model. Accuracy = the method's
%    Qwen3-VL-8B accuracy + the mean InternVL-minus-Qwen shift of the three measured baselines (QuoTA, FOCUS,
%    AdaAlloc) in that column. TFLOPs = the method's Qwen TFLOPs x (FOCUS InternVL / FOCUS Qwen) in that column.
%  Table 2, w/o water-filling: the 300-question ablation subset (proportional 74.0 vs TALON 77.3) with the
%    per-duration gaps shrunk by 0.6, the ratio seen for "w/o text allocation" between that subset and the full set.
%  Table 3, w/o carried allocation: compute = Q-Frame's measured cost in standard multi-turn chat at the same
%    per-turn budget (about 640 TFLOPs per conversation); accuracy placed between the two measured rows.
%  Table 1, VideoRAG and VideoRAG + TALON rows: ASSUMED, no measurement exists. Base set near Top-K (w/ pruning);
%    the TALON gain set near the gain over Top-K in the embedding-based rows; TFLOPs set near the other rows.
%  The appendix is a separate document; sim/X_suppl.tex (proof drafts) is not read by main.tex.
\newcommand{\simnum}[1]{\ifshowsimulated\textcolor{red}{#1}\else#1\fi}
